% arara: clean: {extensions: ['log','aux','synctex.gz','out']}
% arara: pdflatex: { shell: yes }
% arara: pdflatex: { shell: yes }
% arara: clean: {extensions: ['log','aux','synctex.gz','out']}
\documentclass[spanish,addpoints,answers,a4paper]{exam}
\usepackage[T1]{fontenc}
\usepackage[spanish]{babel}
\usepackage{libertine}
\usepackage[tmargin=2cm,rmargin=2cm,lmargin=2cm]{geometry}
\usepackage{amsmath,bm}
\usepackage[shortlabels]{enumitem}
\usepackage{siunitx}
\sisetup{
	group-four-digits = true,
	group-separator = {,}
}
\usepackage{multicol}

\usepackage{graphicx}
\graphicspath{{../../../images/}}
\pagestyle{plain}

\usepackage{minted}
\setminted[cpp]{tabsize=2}
\usepackage{hyperref}
\renewcommand{\solutiontitle}{\noindent\textbf{Solución}\par\noindent}
\renewcommand\listingscaption{Listado}
\newcommand{\unmarkedfntext}[1]{%
	\begingroup
	\renewcommand\thefootnote{}\footnote{#1}%
	\addtocounter{footnote}{-1}%
	\endgroup
}
\spanishdatedel
\decimalpoint

\def\LOGOUNI{%
	\begin{picture}(0,0)\unitlength=1cm
	\put (-13.5,-2.5) {\includegraphics[width=2.8cm]{logouni}}
	\end{picture}
}

\def\LOGOCTIC{%
	\begin{picture}(0,0)\unitlength=1cm
	\put (-8,-1.2) {\includegraphics[height=1cm]{logocticblack}}
	\end{picture}
}

\newcommand{\mychar}{%
	\begingroup\normalfont
	\includegraphics[height=\fontcharht\font`\B]{Octocat.png}%
	\endgroup
}

\begin{document}

\begin{center}
	\sffamily\bfseries\scshape
	{\Large UNIVERSIDAD NACIONAL DE INGENIERÍA}\LOGOUNI\\
	Centro de Tecnologías de la Información y Comunicaciones\LOGOCTIC\\
\end{center}

\vspace{.8cm}

\begin{center}\sffamily\bfseries\large
	Programación en \texttt{C++} \\
	TEMA $\bm{3}$: Estructuras de repetición \mintinline{cpp}{while} y \mintinline{cpp}{for}.
\end{center}

\vspace{.5cm}
\noindent
\makebox[\textwidth]{Nombres y apellidos:\enspace\hrulefill}
\vspace{0.2in}
\makebox[\textwidth]{Nombres y apellidos del instructor:\enspace MSc. Víctor Melchor Espinoza.\hfill}

%\section*{Ejercicios de práctica}
%
%\subsection*{Primera parte}
%
%\begin{questions}
%
%\question Imprimir los números naturales del $1$ al $100$.\label{q:1}
%
%\begin{solution}
%	\begin{listing}[H]
%		\footnotesize
%		\inputminted{cpp}{1a.cc}
%		\caption{Programa \texttt{1a.cc}.}
%		\label{lst:3.1a}
%	\end{listing}
%\end{solution}
%
%\question Modificar el ejercicio~\ref{q:1} para hallar la suma de los $100$ primeros números naturales.
%
%\begin{solution}
%	\begin{listing}[H]
%		\footnotesize
%		\inputminted{cpp}{2a.cc}
%		\caption{Programa \texttt{2a.cc}.}
%		\label{lst:3.2a}
%	\end{listing}
%\end{solution}
%
%\question Hallar la suma de $n$ números ingresados del teclado.
%
%\begin{solution}
%	\begin{listing}[H]
%		\footnotesize
%		\inputminted{cpp}{3a.cc}
%		\caption{Programa \texttt{3a.cc}.}
%		\label{lst:3.3a}
%	\end{listing}
%\end{solution}
%
%\question Diseñe un algoritmo que dada una secuencia de $n$ números enteros diferentes de cero, determinar con la ayuda de un solo bucle:
%\begin{parts}
%\part La cantidad de los números negativos.
%\part La suma de los números negativos.
%\part La cantidad de números positivos.
%\part La suma de los números positivos.
%\end{parts}
%
%\begin{solution}
%	\begin{listing}[H]
%		\footnotesize
%		\inputminted{cpp}{4a.cc}
%		\caption{Programa \texttt{4a.cc}.}
%		\label{lst:3.4a}
%	\end{listing}
%\end{solution}
%
%\question Leer del teclado $20$ números y para cada número imprimir el mismo número, el número elevado al cubo y el número elevado a la cuarta.
%
%\begin{solution}
%	\begin{listing}[H]
%		\footnotesize
%		\inputminted{cpp}{5a.cc}
%		\caption{Programa \texttt{5a.cc}.}
%		\label{lst:3.5a}
%	\end{listing}
%\end{solution}
%
%\question Hallar la suma siguiente:
%\begin{equation*}
%S = 1 + \frac{1}{2} + \frac{1}{3} + \cdots + \frac{1}{n}.
%\end{equation*}
%
%\begin{solution}
%	\begin{listing}[H]
%		\footnotesize
%		\inputminted{cpp}{6a.cc}
%		\caption{Programa \texttt{6a.cc}.}
%		\label{lst:3.6a}
%	\end{listing}
%\end{solution}
%
%\question Hallar la suma de los términos de la siguiente serie
%\begin{equation*}
%S = 1 - \frac{1}{2} + \frac{1}{3} - \frac{1}{4} + \cdots  \pm \frac{1}{n}.
%\end{equation*}
%
%\begin{solution}
%	\begin{listing}[H]
%		\footnotesize
%		\inputminted{cpp}{7a.cc}
%		\caption{Programa \texttt{7a.cc}.}
%		\label{lst:3.7a}
%	\end{listing}
%\end{solution}
%
%\end{questions}
%
%\subsection*{Segunda parte}
%
%\begin{questions}
%
%\question Elabore un programa que presente la siguiente disposición:\label{q:1}
%\begin{equation*}
%	\begin{matrix}
%	11 & 12 & 13 & 14 & 15 \\
%	21 & 22 & 23 & 24 & 25 \\
%	31 & 32 & 33 & 34 & 35
%	\end{matrix}
%\end{equation*}
%
%\begin{solution}
%	\begin{listing}[H]
%		\footnotesize
%		\inputminted{cpp}{1b.cc}
%		\caption{Programa \texttt{1b.cc}.}
%		\label{lst:3.1b}
%	\end{listing}
%\end{solution}
%
%\question A partir del ejercicio~\ref{q:1} obtenga ahora la siguiente disposición:
%\begin{equation*}
%	\begin{matrix}
%	2 & 3 & 4 & 5 & 6 \\
%	3 & 4 & 5 & 6 & 7 \\
%	4 & 5 & 6 & 7 & 8
%	\end{matrix}
%\end{equation*}
%
%\begin{solution}
%	\begin{listing}[H]
%		\footnotesize
%		\inputminted{cpp}{2b.cc}
%		\caption{Programa \texttt{2b.cc}.}
%		\label{lst:3.2b}
%	\end{listing}
%\end{solution}
%
%\question A partir del ejercicio~\ref{q:1} obtenga ahora la siguiente disposición:
%\begin{equation*}
%	\begin{matrix}
%	1 & 1 & 1 & 1 & 1 \\
%	2 & 4 & 8 & 16 & 32 \\
%	3 & 9 & 27 & 81 & 243
%	\end{matrix}
%\end{equation*}
%
%\begin{solution}
%	\begin{listing}[H]
%		\footnotesize
%		\inputminted{cpp}{3b.cc}
%		\caption{Programa \texttt{3b.cc}.}
%		\label{lst:3.3b}
%	\end{listing}
%\end{solution}
%
%\question A partir del ejercicio~\ref{q:1} obtenga ahora la siguiente disposición:
%\begin{equation*}
%	\begin{matrix}
%	- & + & - & + & - \\
%	+ & - & + & - & + \\
%	- & + & - & + & -
%	\end{matrix}
%\end{equation*}
%
%\begin{solution}
%	\begin{listing}[H]
%		\footnotesize
%		\inputminted{cpp}{4b.cc}
%		\caption{Programa \texttt{4b.cc}.}
%		\label{lst:3.4b}
%	\end{listing}
%\end{solution}
%
%\question Diseñe un programa para obtener la siguiente tabla de multiplicar:
%\begin{align*}
%6 &\times 1 &= 6 \\
%6 &\times 2 &= 12 \\
%6 &\times 3 &= 18 \\
%7 &\times 1 &= 7 \\
%7 &\times 2 &= 14 \\
%7 &\times 3 &= 21 \\
%\end{align*}
%
%\begin{solution}
%	\begin{listing}[H]
%		\footnotesize
%		\inputminted{cpp}{5b.cc}
%		\caption{Programa \texttt{5b.cc}.}
%		\label{lst:3.5b}
%	\end{listing}
%\end{solution}
%
%\end{questions}

\section*{Ejercicios obligatorios}

\begin{questions}

\question Elabore un programa que pida al usuario un entero positivo $n$. Luego, el programa debe pedir al usuario $n$ enteros usando un ciclo. El programa deberá encontrar el mayor, menor, y promedio de los valores dados por el usuario e imprimirlos al final.

\begin{solution}
	\begin{listing}[H]
		\footnotesize
		\inputminted{cpp}{1.cpp}
		\caption{Programa \texttt{1.cpp}.}
		\label{lst:1.9}
\end{listing}
\end{solution}

\question Desarrolle un programa que mantenga una cuenta. El programa ofrecerá al usuario un menú de opciones:

\begin{minted}{bash}
0: Limpiar la cuenta.
1: Depositar dinero a la cuenta.
2: Retirar dinero de la cuenta.
3: Mostrar el saldo actual.
4: Salir del programa.
\end{minted}

\begin{solution}
	\begin{listing}[H]
		\footnotesize
		\inputminted{cpp}{2.cpp}
		\caption{Programa \texttt{2.cpp}.}
		\label{lst:1.10}
\end{listing}
\end{solution}

\end{questions}

\begin{flushright}\bfseries
Centro de Tecnologías de la Información y Comunicaciones (CTIC)\\[2mm]
\today\unmarkedfntext{Código disponible en \href{https://github.com/carlosal1015/Cpp-Programming}{\mychar{}}.}
\end{flushright}

\end{document}